%% By design, frex supports arbitrary algebraic theories (including
%% undecidable theories, which can be constructed using quotients,
%% although these introduce proof obligations that the user must
%% discharge by some means other than definitional equality).

\section{Completeness and Certification}
\seclabel{certification}
%
\Frexlib{} uses setoids to extract intensional information: it
automatically extracts printable representations of simplification proofs, and moreover formalises
the completeness of fral/frex simplifiers. Here is how it works:
\begin{itemize}
\item It is easy to construct a frex and a fral using
      a quotient setoid: terms quotiented by provability.
\item The equivalence relation in this setoid is not necessarily
  \emph{effective}/\emph{decidable}.
\item An effective fral/frex must be constructed manually/creatively.
\item Any fral/frex is canonically isomorphic to the quotient frex, and an
  effective fral/frex with a constructive universality proof has an effective isomorphism to
  the quotient setoid.
\item In this way, a simplifier using an effective frex is
  constructively sound and complete.
\item The universal property of the effective fral/frex lets us
  interpret equations it simplifies to the same value as related
  values in every model/extension.
\item
  So the effective fral/frex lets us extract the data structure that represents
  equivalence in the setoid. For the
  quotient fral/frex, it is a representation of the equational proof.
\end{itemize}
Through this mechanism, there is no need to write any proof extraction
code for our simplifiers. The universal properties of the fral/frex
have done all the presentation-specific heavy-lifting. The remainder
of this section expands this description in more technical detail
(\subsecref{completeness}), how \Frexlib{} we represent certificates
(\subsecref{extracting certificates}), and how we simplify proofs
(\subsecref{proof simplification}).

\subsection{Completeness}
\subseclabel{completeness}
%
It is straightforward to represent the deduction rules of equational
logic from \figref{provability} as a datatype
\IdrisType{Provability}$\,t\,s$ relating pairs of terms $t$ and $s$.
There are two variations on this provability relation: for free models
and for free extensions.

For free models over \X{}, we index the relation by $t,s \in \Term
\X$.  Its definition consists of the unshaded rules in
\figref{provability}. The quotient setoid
$Q \definedby \Term\X/\text{\IdrisType{Provability}}$ then satisfies the universal
property of the free model over \X{}.
%
For the free extension of a model $\avar$ by \X{}, we index the
relation by $t,s \in \Term(\X \uplus \U\,\avar)$, i.e., terms with
variables in the disjoint union $\X \uplus \U\,\avar$. The
left-injected values represent variables from \X{}. For each $c \in
\U\,\avar$, the corresponding right-injected variable, which we denote
by $\underline{c}$, is a term representing $c$.  The provability
relation then adds the shaded evaluation equation from
\figref{provability}. It amounts to datatype constructor, for every
operator $f : n$ and constants $c_1, \ldots, c_n \in \U\,\avar$,
representing the equation \( (\underline c_1, \ldots, \underline c_n)
= \underline{\text\avar\sem f(c_1, \ldots, c_n)} \).  The quotient
setoid $Q\definedby \Term(\X \uplus \U\,\avar) /
\text{\IdrisType{Provability}}$ then satisfies the universal property
of the frex of \avar{} by \X.

The interpretation of a term $t$ in $Q$ is $t$ itself, and so the quotient
setoid does not help at all with simplification.  Since $Q$ satisfies
the universal property, it is canonically isomorphic to every other
fral/frex. So the interpretation of every two terms $s,t$ in every
fral/frex is related iff the datatype
$\text{\IdrisType{Provability}}\,t\,s$ is inhabited. Therefore
evaluation in fral/frex is sound and complete for algebraic
simplification.
%
After designing a fral/frex $S$ whose equivalence relation is
\emph{effective}, we can use $Q$ as follows. The input type to the
function \IdrisFunction{solve} is exactly the carrier of $Q$.  By
appealing to the universal property of the fral/frex $S$, we obtain an
inhabitant of \IdrisType{Provability}, and we extracted the
simplification proof, which we call the \emph{certificate}.

\subsection{Representing Certificates}
\subseclabel{extracting certificates} \Frexlib{} provides the
\FrexLemma{} \IdrisBound{pres} type, consisting of a
\IdrisBound{pres}\IdrisFunction{.signature}-equation and derivation
for it in the quotient fral for \IdrisBound{pres}.
Such lemmata are sound: every \FrexLemma{} for a theory holds
in all models of this theory.
%
\Frexlib{} provides a \IdrisFunction{mkLemma} smart constructor which
runs the given fral simplifier, constructs a
proof that a stated equivalence holds, and returns a valid
\IdrisType{Lemma}.
This mechanism allows users to build up a library of
lemmata for their theories. Users can then seamlessly invoke
these lemmata in any model, avoiding further \Frexlib{} calls.

This approach forces the user's project to depend on most of
\Frexlib{} indirectly through such modules.
If users do not want to introduce that dependency, \Frexlib{} gives
them the option not to depend on it.
%
\Frexlib{} also supports proof-certificate extraction,
allowing users to produce standalone libraries of ordinary \idris{} functions
rather than \FrexLemma{} independent of
the \Frexlib{} library. For example, the fral simplifier for monoids
generates a module exporting the following \idris{} function:
\begin{center}
\CertificateType{}
\end{center}
together with an explicit equational proof that does not call any \Frexlib{}
simplification steps (see \figref{certification} in the \appref{printing}).
We now explain the certification process in technical detail.

Our goals for this extraction are to (1) produce libraries from lemmata,
and (2) produce somewhat idiomatic \idris{} code.
%
The derivation found by \Frexlib{} may not be what a human would have
chosen but it should nonetheless be possible for a sufficiently patient
human to follow the reasoning steps.

The main challenge was to go from the rich type of derivations trees,
i.e., the quotient setoid relation, to a representation that we can
print in a relatively readable way. The derivation trees have
arbitrarily nested transitivity, symmetry, and $n$-ary congruence
steps.  We transform them into a type of linear/flat derivations that
could be pretty-printed using combinators for setoid
reasoning.

We represent derivations in layers:
\begin{center}
  \begin{tabular}{@{}l@{\qquad}l@{}}
(a) the reflexive-transitive closure of &
(b) the symmetric closure of\\%
(c) the unary congruence closure of &
(d) axiomatic reasoning steps.
  \end{tabular}
\end{center}
We now detail each of these layers, implementing each layer by a
dedicated data structure~(cf.~\figref{linear-derivations} in
\appref{printing} for the full details):\\
\textit{(a) Reflexive-transitive closure (\IdrisType{RTList}):}
type-aligned~\cite{DBLP:conf/haskell/PloegK14} lists of steps in
the closed-over relation: the target element of each list position
is the source element of the next.\\
\textit{(b) Symmetric closure (\IdrisType{Symmetrise}):}{ either the relation or its opposite}.\\
\textit{(c) Unary congruence closure (\IdrisType{Locate}):} It suffices to pair a term
with a distinguished variable for the contextual hole, together with a
step in the closed-over relation. To ease our pretty-printing code,
we distinguish between using the closed-over relation in an empty
context, and using it in a context with a distinguished variable
represented by the Idris value \IdrisData{Nothing}.\\
\textit{(d) Axiomatic steps (\IdrisType{Step}):} An atomic step is either a setoid equivalence,
or an axiom.\\
Putting these together gives the type \IdrisType{Derivation} of
linear derivations~(cf.~\figref{linear-derivations}(e) in \appref{printing}).

Every derivation tree decomposes into a value in this layered
representation. The modular definition of \IdrisType{Derivation} as a
composition of the relation-transformers \IdrisType{RTList},
\IdrisType{Symmetrise}, \IdrisType{Locate} and \IdrisType{Step} makes
decomposition straightforward.  We use generic combinators for each
closure relation-transformers. Closure under congruence is the
trickiest part, decomposing an $n$-ary congruence in the derivation tree
into $n$ separate
unary congruences, pushing them under the reflexive-transitive and
symmetric closure layers, and erasing any congruence steps with the
identity context.

\subsection{Proof Simplication}
\subseclabel{proof simplification}

Certification also allows us to inspect \Frexlib{}-generated proofs.
%
Frexlet developers can check whether data-structures and
proofs are suboptimal, spurring code refactoring.
%
Concretely, when developing \Frexlib{}, we noticed proofs with loops:
multi-step derivations that start and end in the same term.
%
Such loops come from internal data structures that optimise
simplifier-development effort, but insert semantically irrelevant
subterms that can be simplified away.
%
\Frexlib{} implements a generic proof simplifier that automatically removes
all such loops.
%
This mechanism suggests further work, developing
modules for simplifying these proofs further.
